\documentclass{article}
% \usepackage[utf8]{inputenc} 
\usepackage[T1]{fontenc}   
\usepackage{lmodern}
\usepackage{fullpage}
\usepackage{color, graphicx}
\usepackage{setspace}
\usepackage{epstopdf}
\usepackage{amsmath, amsfonts, amsthm, amssymb, bm, bbm, mathtools, mathrsfs}
\usepackage{dsfont}
\usepackage{scalerel}
\usepackage{verbatim}

\usepackage{array}
\usepackage{booktabs}
\usepackage{multirow}
\usepackage{multicol}
\usepackage{tabu}

\usepackage{algorithm}
\usepackage{algorithmicx}
\usepackage{algpseudocode}

\usepackage[round]{natbib}
\usepackage{xr,xspace}

\usepackage{xcolor}

\newcommand{\red}{\color{red}}
\newcommand{\blue}{\color{blue}}
\newcommand{\green}{\color{green!60!black}}
\definecolor{darkblue}{rgb}{0.0, 0.0, 0.55}
\definecolor{chicago-maroon}{RGB}{128,0,0}

\usepackage[bookmarks, colorlinks=true, plainpages=false, 
            citecolor=darkblue, linkcolor=chicago-maroon, 
            anchorcolor=red, urlcolor=teal]{hyperref}
\usepackage[nameinlink]{cleveref}
\usepackage{prettyref}
            
\usepackage{url}
\urlstyle{rm}

\usepackage{etoolbox}
\usepackage{appendix}
\usepackage{enumitem}
\usepackage{caption,subcaption}
\usepackage{soul}
\usepackage{romanbar}
\usepackage{todonotes}
\usepackage{tcolorbox}
\usepackage{namedtensor}

% \usepackage{almendra}

\usepackage{authblk}

\usepackage{pgf,tikz}

\usepackage{float}
\usetikzlibrary{arrows,arrows.meta,shapes.arrows,shapes.geometric,shapes.multipart,fit,automata,decorations.pathmorphing,positioning,shapes.swigs}
\tikzset{
	%Define standard arrow tip
	>=stealth',
	%Define style for boxes
	true/.style={
		rectangle,
		draw=black, very thick,
		text width=6.5em,
		minimum height=2em,
		text centered,
		fill=gray, opacity = 0.5},
	punkt/.style={
		rectangle,
		rounded corners,
		draw=black, very thick,
		text width=6.5em,
		minimum height=2em,
		text centered},
	est/.style={
		circle,
		draw=black, very thick,
		text centered},
	shade/.style={
		circle,
		draw=black, very thick, fill=gray!50,
		text centered},
	weight/.style={
		circle,
		draw=black, very thick,
		text width=6.5em,
		minimum height=2em,
		text centered},
	% Define arrow style
	pil/.style={
		->,
		thick,
		shorten <=2pt,
		shorten >=2pt,},
	double/.style={
		<->,
		thick,
		shorten <=2pt,
		shorten >=2pt,},
	dash/.style={
		dashed,
		thick,
		shorten <=2pt,
		shorten >=2pt,},
	dashdouble/.style={
		<->,
		dashed,
		thick,
		shorten <=2pt,
		shorten >=2pt,}
}

\usepackage{tikz-cd}
\makeatletter 
\newcolumntype{C}[1]{>{\centering\arraybackslash}p{#1}}

\usepackage{fancyhdr}
\usepackage{microtype} 
\providecommand{\keywords}[1]{\textbf{Keywords:} #1}

\newcommand{\myreferences}{Master.bib}
\newcommand{\myreferencesapp}{Master.bib}

\newcommand{\bbN}{{\mathbb{N}}}

\theoremstyle{plain}
\newtheorem{theorem}{Theorem}
\newtheorem{lemma}{Lemma}
\newtheorem{proposition}{Proposition}
\newtheorem{corollary}{Corollary}
\newtheorem{observation}{Observation}
\theoremstyle{definition}
\newtheorem{definition}{Definition}
\newtheorem{example}{Example}
\newtheorem{hypothesis}{Hypothesis}
\newtheorem{assumption}{Assumption}
\newtheorem{condition}{Condition}
\newtheorem{conjecture}{Conjecture}
\newtheorem{problem}{Problem}
\newtheorem{question}{Question}
\newtheorem{remark}{Remark}
\newtheorem{claim}{Claim}
\newtheorem*{remark*}{Remark}

\newcommand{\bbX}{{\mathbb{X}}}

\newcommand{\calA}{{\mathcal{A}}}
\newcommand{\calB}{{\mathcal{B}}}
\newcommand{\calC}{{\mathcal{C}}}
\newcommand{\calD}{{\mathcal{D}}}
\newcommand{\calE}{{\mathcal{E}}}
\newcommand{\calF}{{\mathcal{F}}}
\newcommand{\calG}{{\mathcal{G}}}
\newcommand{\calH}{{\mathcal{H}}}
\newcommand{\calI}{{\mathcal{I}}}
\newcommand{\calJ}{{\mathcal{J}}}
\newcommand{\calK}{{\mathcal{K}}}
\newcommand{\calL}{{\mathcal{L}}}
\newcommand{\calM}{{\mathcal{M}}}
\newcommand{\calN}{{\mathcal{N}}}
\newcommand{\calO}{{\mathcal{O}}}
\newcommand{\calP}{{\mathcal{P}}}
\newcommand{\calQ}{{\mathcal{Q}}}
\newcommand{\calR}{{\mathcal{R}}}
\newcommand{\calS}{{\mathcal{S}}}
\newcommand{\calT}{{\mathcal{T}}}
\newcommand{\calU}{{\mathcal{U}}}
\newcommand{\calV}{{\mathcal{V}}}
\newcommand{\calW}{{\mathcal{W}}}
\newcommand{\calX}{{\mathcal{X}}}
\newcommand{\calY}{{\mathcal{Y}}}
\newcommand{\calZ}{{\mathcal{Z}}}

\newcommand{\bbR}{{\mathbb{R}}}
\newcommand{\reals}{{\mathbb{R}}}
\newcommand{\naturals}{\mathbb{N}}
\newcommand{\positivenaturals}{\mathbb{N}^+}


% graph notations
\def\allgraphs{\mathcal{G}}
\def\allsymbols{\mathcal{I}}
\newcommand{\subgraphs}[1]{\mathcal{G}[#1]}
\newcommand{\degreevg}[2]{\text{deg}(#1,#2)}

\newcommand{\contractvg}[2]{\mathcal{C}[#1](#2)}
\newcommand{\contractv}[1]{\mathcal{C}[#1]}
\newcommand{\deletevg}[2]{\mathcal{D}[#1](#2)}
\newcommand{\deletev}[1]{\mathcal{D}[#1]}
\newcommand{\degeneracy}[1]{D(#1)}
\newcommand{\degeneracypg}[2]{D[#1](#2)}
\newcommand{\contraction}[1]{C(#1)}
\newcommand{\contractionpg}[2]{C[#1](#2)}

\newcommand{\vertices}[0]{V}
\newcommand{\verticesg}[1]{V(#1)}
\newcommand{\edges}[0]{E}
\newcommand{\edgesg}[1]{E(#1)}
\newcommand{\neighborsvg}[2]{N(#1,#2)}

\newcommand{\completebig}[2]{K_{#1,#2}}
\newcommand{\treewidth}{\mathrm{tw}}
% combination notations
\newcommand{\permutations}[1]{\Pi_{#1}}
\begin{document}

\title{A Graph Family with Treewidth Lower-Bounded by $\omega(\sqrt{|E|})$}
\author{Zhang Ruiqi}
\maketitle

\section{Introduction}
For $n \geq 3$ be an integer, the line graphs of complete bipartite graph $G_n :=L(K_{n,n})$ satisfies
\begin{align*}
    \mathrm{tw}(G_n) = \Theta(|V(G_n)|^{\frac{2}{3}}),
\end{align*}
which implies
\begin{align*}
    \mathrm{tw}(G_n) = \omega(|V(G_n)|^{\frac{1}{2}}).
\end{align*}

This implies that treewidth of a graph $G$ can't be upper bounded by $O(|\edgesg{G}|^{\frac{1}{2}})$.
\section{Preliminaries}
All graphs in this paper are finite, undirected, and without self-loops or multiple edges. 

\begin{definition}[Line Graph]\label{def:line_graph}
Let $G$ be a graph, its line graph $L(G)$ is defined by
\begin{align*}
    \verticesg{L(G)} =& \edgesg{G},\\
    \edgesg{L(G)} =& \{\{u,v\}|u \cap v \neq \emptyset\}.
\end{align*}
\end{definition}

\begin{remark}
    By definition~\ref{def:line_graph}, the vertices of $L(G)$ are edges of the original graph $G$. Vertices $e_1, e_2$ in $\verticesg{L(G)}$ are adjacent in $L(G)$ if and only if they share a common vertex in $G$.
\end{remark}

\begin{remark}
    Following definition~\ref{def:line_graph}, we can easily get the number of edges of $L(G)$.

    \begin{align*}
        |\edgesg{L(G)}| = \sum_{v \in \verticesg{G}} \binom{\degreevg{v}{G}}{2}.
    \end{align*}
    
    Especially for complete bipartite graph $\completebig{n}{n}$, we have
    \begin{align}\label{eq:edges_linegraph_complete_bigraph}
        |\edgesg{L(\completebig{n}{n})}| =& n^2(n-1) \notag \\
        =& \Theta(n^3).
    \end{align}
\end{remark}

For complete bipartite graphs $\completebig{p}{q}$, \cite{HarveyWood2018} establishes a lower bound on the treewidth of their line graphs $L(\completebig{p}{q})$.

\begin{theorem}[Theorem 1.4 in \cite{HarveyWood2018}]\label{thm:tw_lower_bound_of_L(Kpq)}
    For $p,q \in \mathbb{N}$, 
    \begin{align*}
        \mathrm{tw}(L(\completebig{p}{q})) \geq \frac{1}{2}pq-1.
    \end{align*}
\end{theorem}

Especially, if $p= q = n \geq 3$, we can determine $\mathrm{tw}(L(\completebig{n}{n}))$ exactly. This result is given from  Theorem~1.4 in \cite{harvey2014treewidth} and Theorem~3.1 in \cite{lucena2007achievable}. 

\begin{theorem}[Treewidth of $L(\completebig{n}{n})$]\label{thm:tw_L(Knn)}
For any $n \geq 3$ in $\mathbb{N}$, we have
\begin{align*}
    \treewidth(L(\completebig{n}{n})) =& \frac{n^2}{2} + \frac{n}{2} - 1 \\
    =& \Theta(n^2).
\end{align*}
\end{theorem}

\section{Main Result}
The main result of this article is as follows:
\begin{proposition}\label{pro:unbounded_graphs}
    There exists a graph sequence $\{G_n\}$, such that
    \begin{align*}
        \treewidth(G_n) = \omega(|\edgesg{G_n}|^{\frac{1}{2}}.
    \end{align*}
\end{proposition}
\begin{proof}
    Let $G_n = L(\completebig{n+2}{n+2}), n \in \mathbb{N}_+$.

    Following Theorem~\ref{thm:tw_L(Knn)}, we have
    \begin{align}\label{eq:tw_unbounded_graph}
        \treewidth(G_n) = \Theta(n^2).
    \end{align}

    Following Equation~\eqref{eq:edges_linegraph_complete_bigraph}, we have
    \begin{align}\label{eq:number_edges_unbounded_graph}
        |\edgesg{G_n}| = \Theta(n^3).
    \end{align}

    Collecting equation~\eqref{eq:tw_unbounded_graph} and equation~\eqref{eq:number_edges_unbounded_graph}, we have
    \begin{align*}
        \treewidth(G_n) =& \Theta({|\edgesg{G_n}|^\frac{2}{3}})\\
        =& \omega({|\edgesg{G_n}|^\frac{1}{2}}).
    \end{align*}
    which completes the proof.
\end{proof}

\begin{remark}
    Each vertex of $L(\completebig{n}{n})$ is of degree $2n-2$, which is even. This implies that $L(\completebig{n}{n})$ is an Euler graph, and must be corresponding to a V-statistic we need when computing $n^2(n-1)$-order U-statistic of our form.

    For example, $L(\completebig{3}{3})$ is corresponding to a V-statistic with mode
    \begin{align*}
        (ab, bc, cd, de, ef, fd, db, bg, gh, hc, ci, ie, ea, ag, gf, fh, hi, ia)
    \end{align*}
    which is need to be computed when computing $18$-order U-statistic.
\end{remark}
\bibliographystyle{plainnat}
\bibliography{Master}
\end{document}

